\subsection{Task 1-4: DCOPF With Line Constraints}
% ----------------------------------------------------------

The thermal limits are
$P_{12}^{\max} = 500$~MW, $P_{13}^{\max} = 500$~MW,
$P_{23}^{\max} = 100$~MW.

% ---------- 1-4a ----------
\subsubsection*{(a) Qualitative effect of line constraints}

Without line constraints, Generator~1 (cheapest at 300~NOK/MWh) would supply
the entire 900~MW load, leading to $P_{23} \approx 218$~MW
(see Task~1-3 flows).  This violates the 100~MW limit on Line~2-3.

The binding constraint on Line~2-3 prevents the unconstrained-optimal
power pattern from being realised.  As a result, Generator~3 at Node~3 must
pick up part of the local load rather than relying on cheap remote power, and
the system marginal price is no longer uniform: nodes separated by the
congested line will see different nodal prices.  Lines 1-2 and 1-3 are not
binding at the constrained optimum.

% ---------- 1-4b ----------
\subsubsection*{(b) Additional term in the Lagrangian}

For each transmission line $ij$ with an active upper flow limit, a
non-negative multiplier $\mu_{ij} \geq 0$ is added.  With only Line~2-3
binding:

\[
  \mathcal{L} = \sum_{n=1}^{3} c_n P_{Gn}
              + \lambda\!\left(\sum_{n} D_n - \sum_{n} P_{Gn}\right)
              + \mu_{23}\!\left(P_{23} - P_{23}^{\max}\right)
\]

where $\lambda$ is the system energy price (dual of the power balance) and
$\mu_{23} \geq 0$ is the shadow price (congestion rent) of Line~2-3.
The DC flow is $P_{23} = b_{23}(\delta_2 - \delta_3)$, which can be
expressed directly in terms of generator outputs via the B-matrix
(see below).

% ---------- 1-4c ----------
\subsubsection*{(c) First-order (KKT) conditions}

Using the DC power flow, the line flow can be written as a linear function
of net nodal injections.  The sensitivity of $P_{23}$ (in MW) to each
generator output is obtained from the reduced B-matrix
$\mathbf{B}_{\text{red}} = \bigl[\begin{smallmatrix}50&-30\\-30&40\end{smallmatrix}\bigr]$
(buses 2 and 3, with bus~1 as slack, susceptances in per unit,
$S_{\text{base}}=100$~MVA):

\[
  \frac{\partial P_{23}}{\partial P_{G1}} = 0,\qquad
  \frac{\partial P_{23}}{\partial P_{G2}} = \frac{3}{11},\qquad
  \frac{\partial P_{23}}{\partial P_{G3}} = -\frac{6}{11}
\]

The KKT stationarity conditions are:

\[
  \frac{\partial \mathcal{L}}{\partial P_{G1}} = 0
  \;\Rightarrow\; c_1 - \lambda = 0
  \;\Rightarrow\; \lambda = c_1 = 300~\text{NOK/MWh}
\]

\[
  \frac{\partial \mathcal{L}}{\partial P_{G3}} = 0
  \;\Rightarrow\; c_3 - \lambda - \frac{6}{11}\,\mu_{23} = 0
  \;\Rightarrow\; \mu_{23} = \frac{11(c_3 - \lambda)}{6}
\]

\[
  \frac{\partial \mathcal{L}}{\partial P_{G2}} \geq 0
  \;\text{(with } P_{G2} = 0\text{)}
  \;\Rightarrow\; c_2 - \lambda + \frac{3}{11}\,\mu_{23} \geq 0
\]

The nodal price (LMP) at each bus is
$\lambda_n = \lambda - \mu_{23}\cdot\mathrm{PTDF}_{23,n}$,
where $\mathrm{PTDF}_{23,n}$ is the power transfer distribution factor for
Line~2-3 with injection at bus~$n$ and withdrawal at the slack bus~1:

\[
  \mathrm{PTDF}_{23,1} = 0,\qquad
  \mathrm{PTDF}_{23,2} = \frac{3}{11},\qquad
  \mathrm{PTDF}_{23,3} = -\frac{6}{11}
\]

% ---------- 1-4d ----------
\subsubsection*{(d) Numerical solution}

\paragraph{Optimal generation.}

Without any line limit, $P_{G2} = P_{G3} = 0$ and $P_{23} \approx 218$~MW,
which violates $P_{23}^{\max} = 100$~MW.  The constraint on Line~2-3 is
therefore binding at the optimum.

Since Generator~2 is the most expensive unit (1000~NOK/MWh), the cost is
minimised by setting $P_{G2} = 0$ and distributing the remaining generation
between $G_1$ (300~NOK/MWh) and $G_3$ (600~NOK/MWh).

With $P_{G2} = 0$, the DC flow on Line~2-3 expressed via the B-matrix is

\[
  P_{23} = \frac{3\bigl[(P_{G2} - D_2) - 2(P_{G3} - D_3)\bigr]}{11}
         = \frac{3(800 - 2\,P_{G3})}{11}
\]

Setting $P_{23} = 100$~MW and solving:

\[
  \frac{3(800 - 2\,P_{G3})}{11} = 100
  \;\Rightarrow\; 800 - 2\,P_{G3} = \frac{1100}{3}
  \;\Rightarrow\; P_{G3} = \frac{650}{3} \approx 216.7~\text{MW}
\]

\[
  P_{G1} = 900 - P_{G3} = \frac{2050}{3} \approx 683.3~\text{MW}
\]

\[
  \boxed{P_{G1} = \tfrac{2050}{3} \approx 683.3~\text{MW},\quad
         P_{G2} = 0~\text{MW},\quad
         P_{G3} = \tfrac{650}{3} \approx 216.7~\text{MW}}
\]

\paragraph{Total system cost.}

\[
  TC = \frac{2050}{3}\times 300 + 0 + \frac{650}{3}\times 600
     = 205\,000 + 130\,000
     = \mathbf{335\,000~\text{NOK}}
\]

\paragraph{Phase angles.}

The net nodal injections in per unit are

\[
  P_1 = \frac{2050/3 - 200}{100} = \frac{29}{6}~\text{p.u.},\quad
  P_2 = \frac{0 - 200}{100} = -2~\text{p.u.},\quad
  P_3 = \frac{650/3 - 500}{100} = -\frac{17}{6}~\text{p.u.}
\]

With $\delta_1 = 0$, the nodal equations are:

\[
  50\,\delta_2 - 30\,\delta_3 = -2
  \qquad\text{and}\qquad
  -30\,\delta_2 + 40\,\delta_3 = -\tfrac{17}{6}
\]

From the first equation: $\delta_2 = -0.04 + 0.6\,\delta_3$.
Substituting into the second:

\[
  -30(-0.04 + 0.6\,\delta_3) + 40\,\delta_3 = -\tfrac{17}{6}
  \;\Rightarrow\; 1.2 - 18\,\delta_3 + 40\,\delta_3 = -\tfrac{17}{6}
  \;\Rightarrow\; 22\,\delta_3 = -\tfrac{121}{30}
\]

\[
  \delta_3 = -\frac{11}{60} \approx -0.183~\text{rad},\qquad
  \delta_2 = -\frac{3}{20} = -0.150~\text{rad}
\]

\[
  \delta_1 = 0,\qquad
  \delta_2 = -0.150~\text{rad},\qquad
  \delta_3 \approx -0.183~\text{rad}
\]

\paragraph{Power flows.}

\[
  P_{12} = b_{12}(\delta_1 - \delta_2)
         = 20\!\left(0 - \bigl(-\tfrac{3}{20}\bigr)\right)
         = 3~\text{p.u.}
         \;\Rightarrow\; \mathbf{300~\text{MW (bus 1 $\to$ bus 2)}}
\]

\[
  P_{13} = b_{13}(\delta_1 - \delta_3)
         = 10\!\left(0 - \bigl(-\tfrac{11}{60}\bigr)\right)
         = \tfrac{11}{6}~\text{p.u.}
         \;\Rightarrow\; \mathbf{\tfrac{550}{3} \approx 183.3~\text{MW (bus 1 $\to$ bus 3)}}
\]

\[
  P_{23} = b_{23}(\delta_2 - \delta_3)
         = 30\!\left(-\tfrac{3}{20} - \bigl(-\tfrac{11}{60}\bigr)\right)
         = 30 \times \tfrac{1}{30}
         = 1~\text{p.u.}
         \;\Rightarrow\; \mathbf{100~\text{MW (bus 2 $\to$ bus 3)}} \;\checkmark
\]

\noindent\textbf{Nodal balance check:}
\begin{align*}
  \text{Bus 1:} &\quad P_{12} + P_{13} = 300 + \tfrac{550}{3} = \tfrac{1450}{3}~\text{MW}
                 = P_{G1} - D_1 \checkmark \\
  \text{Bus 2:} &\quad -P_{12} + P_{23} = -300 + 100 = -200~\text{MW}
                 = P_{G2} - D_2 \checkmark \\
  \text{Bus 3:} &\quad -P_{13} - P_{23} = -\tfrac{550}{3} - 100 = -\tfrac{850}{3}~\text{MW}
                 = P_{G3} - D_3 \checkmark
\end{align*}

\paragraph{Shadow price of Line~2-3.}

From the KKT condition for Generator~3:

\[
  \mu_{23} = \frac{11(c_3 - \lambda)}{6}
           = \frac{11 \times (600 - 300)}{6}
           = \mathbf{550~\text{NOK/MWh}}
\]

\paragraph{Nodal prices.}

\[
  \lambda_1 = \lambda - \mu_{23} \cdot 0 = 300~\text{NOK/MWh}
\]

\[
  \lambda_2 = \lambda - \mu_{23} \cdot \tfrac{3}{11}
            = 300 - 550 \times \tfrac{3}{11}
            = 300 - 150
            = \mathbf{150~\text{NOK/MWh}}
\]

\[
  \lambda_3 = \lambda - \mu_{23} \cdot \bigl(-\tfrac{6}{11}\bigr)
            = 300 + 550 \times \tfrac{6}{11}
            = 300 + 300
            = \mathbf{600~\text{NOK/MWh}}
\]

% ---------- 1-4e ----------
\subsubsection*{(e) Reasoning for nodal prices, in particular Node~2}

The nodal price at bus~2 is $\lambda_2 = 150$~NOK/MWh, which is \emph{lower}
than the system energy price $\lambda = 300$~NOK/MWh.  Three factors explain
this:

\begin{enumerate}
  \item \textbf{Bus~2 is on the exporting side of the congested line.}
        Line~2-3 carries 100~MW \emph{away} from bus~2 toward bus~3.  An
        additional MW of demand at bus~2 reduces the net export from bus~2,
        which \emph{relaxes} the binding flow constraint on Line~2-3.

  \item \textbf{Relaxing the constraint allows a cheaper dispatch.}
        With the constraint slightly relaxed, the system can substitute
        expensive Generator~3 (600~NOK/MWh) with cheaper Generator~1
        (300~NOK/MWh), lowering total cost.

  \item \textbf{The net effect is a below-system-price LMP at bus~2.}
        A 1~MW increase in demand at bus~2 causes Generator~1 to increase
        by 1.5~MW and Generator~3 to decrease by 0.5~MW, giving
        $\Delta TC = 300\times1.5 + 600\times(-0.5) = 150$~NOK/MWh $= \lambda_2$.
        Because the counter-flow benefit (saved G3 cost) partly offsets the
        extra G1 cost, the LMP at bus~2 is below $\lambda_1$.
\end{enumerate}

% ---------- 1-4f ----------
\subsubsection*{(f) Effect of increasing demand at Node~2 to 201~MW}

The marginal cost of one additional MWh at Node~2 is the nodal price there:

\[
  \Delta C = \lambda_2 \cdot \Delta P_{D2} = 150 \times 1 = \mathbf{150~\text{NOK}}
\]

The new optimal dispatch (with $P_{G2} = 0$ still optimal) is found from

\[
  \frac{3(799 - 2\,P_{G3})}{11} = 100
  \;\Rightarrow\; P_{G3} = \frac{1297}{6} \approx 216.2~\text{MW},\quad
  P_{G1} = \frac{4109}{6} \approx 684.8~\text{MW}
\]

\[
  TC_{\text{new}}
    = \tfrac{4109}{6}\times300 + \tfrac{1297}{6}\times600
    = 205\,450 + 129\,700
    = 335\,150~\text{NOK}
\]

\[
  \Delta TC = 335\,150 - 335\,000 = \mathbf{150~\text{NOK}}
\]

This confirms that $\lambda_2 = 150$~NOK/MWh is indeed the marginal system
cost of an extra MW of demand at Node~2.
